(a) Os dados contêm os dois tipos classicos de erro.

\begin{itemize}
\item \textbf{Erro sistematico:} se os medidores analogicos estiverem descalibrados, todas as leituras de $V$ ou de $I$ podem ficar deslocadas para cima ou para baixo.
\item \textbf{Erro aleatorio:} ha espalhamento visivel nos pontos, causado por resolucao limitada, leitura do ponteiro e pequenas flutuacoes experimentais.
\end{itemize}

(b) Colocando $I$ no eixo horizontal e $V$ no eixo vertical, os pontos
\[(0.8,10),\ (1.1,20),\ (2.5,30),\ (4.2,40),\ (4.3,50),\ (4.7,60),\ (5.8,70),\ (6.4,80)\]
ficam aproximadamente alinhados. Uma reta "a olho" razoavel e
\[
V\approx 12I,
\]
ou, se quisermos admitir pequeno deslocamento vertical,
\[
V\approx 11.5I+2.
\]

Em ambos os casos, a inclinacao sugere resistencia da ordem de $12\,\Omega$, com os pontos de baixa corrente e os proximos de $I\approx4.2$--$4.3$ A mostrando o maior espalhamento em torno da reta.